% =============================================================================
% 6.2 研究局限性
% 草稿版本：v0.2 (2026-05-07)
%   v0.1：作为原 5.4.1 \subsubsection*{局限性回顾} 起草。
%   v0.2：按导师意见 3 拆出，提级为 6.2 \subsection。叙述上把"将在
%        ref{sec:future-work} 节作为未来工作展开"改为"对应 6.3 节"，
%        与新结构对齐。
% 字数：约 700 字
% 风格规则：v0.2 风格（无 \textbf / \textit / 引例括号 / 双引号 / 破折号）
% 引用：\cite{ref23, ref29}
% 图表：无
% 依赖宏包：booktabs / tabularx（主稿已加载）
% =============================================================================

\subsection{研究局限性}
\label{sec:limitations}

本文工作在毕业设计研究范围内对气象 Agent 单点架构、双通路 RAG 机制、
代码生成与受限沙箱、报告生成方案进行了系统化设计与实证。\ref{chap:evaluation}
章的 5 组评测在揭示有效性的同时也暴露了若干局限性，本节按局限性来源
归纳为四类，每一类都自然指向 \ref{sec:future-work} 节的对应未来工作。

第一类局限性是多轮长程依赖能力不足，源自 \ref{sec:intent-eval} 节
意图识别评测。\texttt{context\_\bk{}inference} 类 3 条用例端到端通过率仅
为 33.3\%，与单轮场景的 100\% 之间形成 66.7 个百分点差距。当查询
同时切换意图与隐式继承多个槽位时，大语言模型倾向于把所有未在 query
中显式出现的字段都视为重置而非继承。这一现象表明单纯把上下文摘要
拼接到 prompt 不足以实现稳定的多轮槽位继承，需要在工程上引入
显式的对话状态跟踪机制并与意图识别管线协同。

第二类局限性是数值阈值的精确归属仍依赖确定性分级器兜底，源自
\ref{sec:rag-eval} 节通路 A 检索评测。即使在最优混合权重 $\alpha=0.9$
下，60 条评测集上仍有约 9 条 Top-1 脱靶，且全部集中在邻近阈值边界、
术语别名冲突、跨类别多文档三类场景。混合检索的工程上限在 85\% 左右，
难以单纯依靠权重调整突破。本文通过通路 B 即确定性 if-else 分级器
加 \texttt{grade\_\bk{}id} 硬链接的方式覆盖了这一盲区，但本质上是把问题
绕过而非解决。如果未来要扩展到行业气象等复杂领域，分级器规则的
人工编写成本会显著上升。

第三类局限性是受限沙箱的安全性是轻量原型而非生产级方案。
\ref{sec:code-eval} 节的沙箱仅做了 builtins 与模块的白名单约束加
墙钟超时加结构化错误捕获四项，能够阻止意外死循环与文件读写，但不
能抵御主动恶意攻击。生产环境下应当使用 RestrictedPython、Pyodide
或 Docker 容器实现进程级隔离。本文限定在毕业设计研究
范围内主动放弃了这一安全维度，保留为未来工作。

第四类局限性是集成衰减效应，源自 \ref{sec:e2e-eval} 节端到端评测。
通路 B 单测 \texttt{citation} 真实性已达 100\%，但端到端集成后衰减
到 26.7\%，意味着大语言模型在最终报告生成阶段仍可能省略检索到的
GB/T 编号。这是一类只有在集成层评测中才会暴露的问题，单独优化任何
一个模块都无法消除。

上述四类局限性既是本文研究工作的边界，也构成了下一节未来工作的
直接动因。每一类局限性都对应一个或多个具体的工程改进方向，
\ref{sec:future-work} 节按这一对应关系展开 4 类未来工作。
